% !TeX root = ../presentacion.tex
% ===========================================================================
%  11. Cierre — 0:40
%
%  Cierra el círculo con los seis criterios de aceptación del documento de
%  alcance (§7). Es la tabla que el evaluador busca, y la misma que el
%  informe (E1, §7.2) sostiene con evidencia.
%
%  REVISAR LA VÍSPERA, CON EL SISTEMA DELANTE: marcar con \si solo lo que se
%  haya visto correr ese día. Un criterio dado por bueno que luego falla en
%  la demo cuesta más que uno declarado pendiente.
% ===========================================================================

\bloque{Cierre}{0:40}{\exponeArquitectura}{Todos}

\begin{frame}{Los seis criterios de aceptación}
  \note{Leer solo la columna de la derecha, y no adornar. Si algo está a
  medias, se dice.}

  \begin{center}
    \scriptsize
    \setlength{\tabcolsep}{3pt}
    \renewcommand{\arraystretch}{1.05}
    \begin{tabular}{@{}cL{4.2cm}L{5.1cm}@{}}
      \toprule
      & \textbf{Criterio del alcance (§7)} & \textbf{Evidencia}\\
      \midrule
      \si & \crit{1} El socio consulta, reserva y cancela & Demo, pasos 1--5; \id{BookingServiceCreacionTest}\\
      \si & \crit{2} El admin gestiona canchas y cancela cualquier reserva & Demo, pasos 6--7; \id{CourtServiceTest}\\
      \si & \crit{3} No se reserva un bloque ocupado & Demo, paso 3; \rn{02} y \id{ConcurrenciaReservaIT}\\
      \si & \crit{4} Reportes de ocupación y reservas por período & Demo, paso 8; los cuatro indicadores\\
      \si & \crit{5} Shell y $\geq$ 2 remotes con Module Federation & Tres remotes, desplegables por separado\\
      \si & \crit{6} $\geq$ 2 microservicios sobre PostgreSQL & Cuatro servicios, tres bases aisladas\\
      \bottomrule
    \end{tabular}
  \end{center}

  \vspace{0pt}
  {\scriptsize\textbf{Lo que nos llevamos:} una regla crítica que solo vive en
  el código de aplicación no es una regla, es una intención; y la
  independencia de datos hay que hacerla imposible de violar.}
\end{frame}

\begin{frame}[standout]
  Gracias\\[6pt]
  {\normalsize ¿Preguntas?}
\end{frame}
